\section{Introduction}\label{ap:introduction}

Erd\H{o}s asked whether every set $A$ of positive integers for which
$\sum_{a\in A}1/a$ diverges contains arithmetic progressions of every
finite length~\cite[Problem~4.33.6]{Erdos1974}.
We prove a quantitative density bound that implies this assertion.
For integers $k\ge3$ and $N\ge1$, let $r_k(N)$ be the largest
cardinality of a subset of $[N]=\{1,\ldots,N\}$ containing no
nonconstant $k$-term arithmetic progression
$a,a+d,\ldots,a+(k-1)d$, with $d>0$.
All logarithms are natural.

\begin{theorem}\label{ap:main}
For each fixed integer $k\ge3$ there are constants
$C_k,c_k,\varepsilon_k>0$ such that, for every $N\ge2$,
\begin{equation}\label{ap:density-bound}
 r_k(N)\le C_kN\exp\bigl(-c_k(\log N)^{\varepsilon_k}\bigr).
\end{equation}
Equivalently, there is $A_k\ge1$ such that an $\alpha$-dense subset
of $[N]$ contains a nonconstant $k$-term progression whenever
\begin{equation}\label{ap:threshold}
 \log N\ge A_k\bigl(2+\log(1/\alpha)\bigr)^{A_k},
 \qquad 0<\alpha\le1.
\end{equation}
\end{theorem}


The constants may depend on the fixed length $k$; the exponent is
not optimized. The reciprocal-sum consequence follows immediately
from summing the density bound over dyadic intervals.

\begin{corollary}[Divergent reciprocal sums]\label{ap:reciprocal}
Every set $A\subseteq\mathbb N$ with
$\sum_{a\in A}a^{-1}=\infty$ contains nonconstant arithmetic
progressions of every finite length.
\end{corollary}

\begin{proof}
Fix $k\ge3$, and suppose that $A$ contains no nonconstant
$k$-term progression. For $m\ge0$, translation identifies
$[2^m,2^{m+1})\cap\mathbb Z$ with an interval of length $2^m$.
Thus $|A\cap[2^m,2^{m+1})|\le r_k(2^m)$. Theorem~\ref{ap:main}
implies, for $m\ge1$,
\[
 \sum_{a\in A\cap[2^m,2^{m+1})}\frac1a
 \le 2^{-m}r_k(2^m)
 \le C_k\exp\bigl(-c_k(m\log2)^{\varepsilon_k}\bigr).
\]
For every $d,\varepsilon>0$, the sequence $e^{-dm^\varepsilon}$
is summable: since $m^\varepsilon/\log m\to\infty$, its terms
are at most $m^{-2}$ for all sufficiently large $m$.
The $m=0$ interval contributes at most one. The reciprocal sum of
$A$ is therefore finite, a contradiction. Divergence also implies
that $A$ is infinite, which supplies lengths one and two.
\end{proof}

\subsection{Density thresholds and summability}

The threshold in \eqref{ap:threshold} is quasipolynomial in
$1/\alpha$: the logarithm of a sufficient interval length is a fixed
power of $2+\log(1/\alpha)$. The inversion between the two forms of
Theorem~\ref{ap:main}, including the endpoints, is proved in
Section~\ref{ap:iteration}.

The infinite-set consequence depends on more than a density tending
to zero. Divergence of a reciprocal sum allows zero-density sets,
and the qualitative estimate $r_k(N)=o(N)$ does not by itself settle
Erd\H{o}s's question. The proof of Corollary~\ref{ap:reciprocal}
uses the stronger condition
\[
                   \sum_{m\ge1}\frac{r_k(2^m)}{2^m}<\infty.
\]
Our bound also gives a stronger divergence criterion: for every
fixed $B\ge0$, a set $A\subseteq\mathbb N$ satisfying
\[
             \sum_{a\in A}\frac{(\log(2+a))^B}{a}=\infty
\]
contains progressions of every finite length. Section~\ref{ap:weighted}
proves this criterion, uniform harmonic bounds and tails, and
estimates for larger weights whose parameters depend on $k$.
It also recovers the theorem of Green and Tao that every subset
of the primes of positive relative upper density contains infinitely
many progressions of every finite length~\cite[Theorem~1.2]{GreenTaoPrimes2008}.

\subsection{Prior density bounds}

Van der Waerden's theorem concerns a partition: every finite coloring
of the positive integers has a monochromatic arithmetic progression
of each finite length~\cite{vanDerWaerden1927}. The density problem
asks for a condition on a prescribed set. Erd\H{o}s and Tur\'an
began the study of the corresponding extremal functions in 1936
and explicitly conjectured that $r_3(N)=o(N)$~\cite{ErdosTuran1936}.
The all-length positive-density conjecture appears explicitly in
Erd\H{o}s's later account~\cite[Section~I.19]{Erdos1961}.
Roth proved the three-term case by Fourier-analytic
methods~\cite{Roth1953}; Szemer\'edi proved the four-term case in
1969 and the general theorem in 1975~\cite{Szemeredi1969,Szemeredi1975}.
Furstenberg's subsequent proof recast the assertion as multiple
recurrence in a measure-preserving system~\cite{Furstenberg1977}.
The reciprocal-sum conjecture extends this density problem to
sparser sets, including the primes~\cite[Problem~4.33.6]{Erdos1974}.

For three terms, successive quantitative advances sharpened Roth's
density-increment method: a shortage of progressions forces a density
increase on a smaller structured set. Roth obtained
$r_3(N)\ll N/\log\log N$. Heath-Brown and Szemer\'edi replaced
the double logarithm by a positive power of $\log N$
\cite{HeathBrown1987,Szemeredi1990}, and Bourgain developed
density increments on regular Bohr sets~\cite{Bourgain1999,Bourgain2008}.
Sanders reached $r_3(N)\ll N/(\log N)^{1-o(1)}$, and Bloom
improved the remaining logarithmic factors~\cite{Sanders2011,Bloom2016}.
Bloom and Sisask gave an almost-periodicity approach to this
near-logarithmic range~\cite{BloomSisask2019}. Their later bound
$r_3(N)\ll N/(\log N)^{1+c}$, for an absolute $c>0$,
crossed the dyadic summability threshold and established the
three-term reciprocal-sum conjecture~\cite{BloomSisask2020}.

Kelley and Meka obtained the stronger bound
\[
             r_3(N)\ll N\exp\bigl(-c(\log N)^{1/12}\bigr)
\]
by combining positivity, dependent-random-choice sifting and
almost-periodicity to control the loss in a density
increment~\cite{KelleyMeka2023}. Bloom and Sisask simplified the
argument and then improved the exponent to
$1/9$~\cite{BloomSisaskKelleyMeka,BloomSisaskImprovement2023}.
Raghavan subsequently obtained
$r_3(N)\le N\exp\bigl(-(\log N)^{1/6-o(1)}\bigr)$
\cite[Theorem~1.4]{Raghavan2026}. These refinements concern three
terms. On the lower-bound side, Behrend's construction already gives
$r_3(N)\ge N\exp(-C\sqrt{\log N})$, so a fixed power saving
in $N$ is impossible even in this case~\cite{Behrend1946}.

Gowers's approach to longer progressions detects higher-order
structure. He introduced the uniformity norms that control
progression counts and extended density increments beyond linear
Fourier analysis~\cite{Gowers1998,Gowers2001}. His general theorem
gives $r_k(N)\ll_k N/(\log\log N)^{c_k}$ for every fixed $k$,
with $c_k>0$. Green and Tao obtained
$r_4(N)\ll N/(\log N)^c$~\cite{GreenTao2017}.
Leng, Sah and Sawhney first treated five terms and then proved,
for every fixed $k\ge5$,
\[
 r_k(N)\ll_k N\exp\bigl(-(\log\log N)^{c_k}\bigr),
 \qquad 0<c_k<1
\]
\cite{LengSahSawhneyFive,LSSSzemeredi2024}.
Their general argument combines a quasipolynomial inverse theorem
with density increments. The inverse theorem controls the structured
correlation obtained from a large uniformity norm; the cost of
localizing and iterating that correlation remains a separate issue.
The resulting saving in $\log\log N$ does not give the dyadic
summability above.

Theorem~\ref{ap:main} gives a stretched-exponential saving in
$\log N$ for every fixed length. Its contribution to Erd\H{o}s's
conjecture is this all-length, summable bound; no improvement of
the three-term exponent is claimed. We next explain the
quantitative constraint on the density-increment regions that
allows the argument to be iterated at this scale.

\subsection{Preserving structure through density increments}

A fixed multiplicative density gain can be applied only
$O_k(1+\log(1/\alpha))$ times before the density exceeds one.
To obtain Theorem~\ref{ap:main}, the total logarithmic loss of
interval length over these increments must be polynomial in
$2+\log(1/\alpha)$. We therefore need to control both the number
of coordinates describing each new region and the precision of
the constraints inherited from previous regions.

The regions are triangular polynomial cells. A spatial integer
variable $u$ determines integer blocks $b_1,\ldots,b_D$
successively through constraints of the form
\[
 \|b_h-C_h(u,b_1,\ldots,b_{h-1})-l_h\|_\infty\le w_h,
 \qquad 1\le h\le D,\qquad 0<w_h\le1/32.
\]
Here $D$ depends only on $k$, each $l_h$ is a fixed center, and
$C_h$ is a polynomial. The spatial variable has weight one,
the block $b_h$ has weight $h$, and $C_h$ has weighted degree
at most $h$. The narrow widths make each integer block unique
when it exists. Keeping these integer blocks, and substituting
them exactly into all later polynomials, preserves the equations
defining the old region when a new constraint is added.
No bound on the coefficients of the $C_h$ is imposed.

Write $Q_h=\log(2/w_h)$ for the precision at weight $h$.
The main increment theorem bounds the \emph{additional}
precision at weight $h$ by a polynomial in the dimension data,
the logarithmic density parameter, and the precisions at
strictly higher weights. In particular, this bound does not use
the current $Q_h$ or the lower precisions. The dimensions obey
an analogous descending recurrence. Since the number $D$ of
weights is fixed, descending induction bounds all dimensions
and precisions by polynomials through the required number of
increments. Section~\ref{ap:setup} states these recurrences
precisely; Section~\ref{ap:iteration} proves that the construction
satisfies them.

Two different arguments supply these bounds. The first produces
an \emph{absolute increment} for a $[0,1]$-valued function on an
integer box with sufficiently long sides. The function must have
positive mean and contain no
nonconstant $k$-term progression in its support. The increment
adds only a number of
determining coordinates polynomial in the logarithmic inverse
mean, independently of the dimension of that box. The second
returns such an increment from a sampled box to the original
cell, preserving the old determining equations and the
descending dependence of their width losses. An absolute
increment with uncontrolled dependence on the old widths would
not suffice for this second task.

\paragraph{The absolute increment.}
The higher-order inverse theorem converts a large uniformity
norm into correlation with a nilsequence, a Lipschitz function
evaluated along a polynomial orbit on a compact nilmanifold.
This description of structured correlations has both ergodic
and additive-combinatorial origins. Host and Kra developed
seminorms and nilsystem characteristic factors, Ziegler proved
the universal-factor theorem, and Bergelson, Host and Kra
isolated nilsequences in multiple correlations
\cite{HostKra2005,Ziegler2007,BergelsonHostKra2005}.
Green and Tao obtained locally quadratic structure for the
$U^3$ norm~\cite{GreenTaoU32008}; Green, Tao and Ziegler then
proved the all-degree interval inverse
theorem~\cite{GreenTaoZiegler2012}.
Manners gave explicit quantitative inverse bounds on prime
cyclic groups~\cite{Manners2018}. The stronger quasipolynomial
bounds of Leng, Sah and Sawhney are the interval input used
here~\cite[Theorem~1.2]{LSSInverse2024}; the appendices prove
the box and product-cyclic reductions needed in the argument.

To pass from structured tests to positive progression counting,
we prove a shift-comparison theorem: lowering the testing degree
permits multiplication by a translated nonnegative function
outside a small exceptional set of shifts. Its degree-one
argument adapts Kelley--Meka unbalancing and sifting, as
developed on Bohr sets by Bloom and Sisask
\cite{KelleyMeka2023,BloomSisaskKelleyMeka}.
The almost-periodicity input is Schoen--Sisask's
Theorem~5.4~\cite{SchoenSisask2016}. It combines the probabilistic
method of Croot and Sisask with Sanders's local control of large
spectra~\cite{CrootSisask2010,SandersSumsets2008}; see also
Sanders's Bogolyubov--Ruzsa theorem~\cite{SandersBogolyubov2012}.
The Schoen--Sisask radius bound retains the incoming radius as a multiplicative
factor, which is essential for the shift comparison.

Higher-degree comparison requires control of polynomial
orbits. Leibman established a qualitative distribution
theorem~\cite{Leibman2005}. Green and Tao made this theory
quantitative and factored polynomial orbits into smooth,
equidistributed and rational components~\cite{GreenTaoOrbits2012}.
Leng's efficient estimates reduce the nilpotency step rather
than repeatedly reducing the dimension
\cite{LengEquidistribution2023}.
We use this mechanism, together with the inverse theorem,
to remove top-degree frequencies in the comparison argument.
The appendices prove the required forms with explicit dependence
on all growing dimensions and parameter counts.

The comparison and counting arguments first give the absolute
rule in a fixed dimension. A relative lifting theorem extends
it with only an additive number of new determining coordinates.
Applied with the constant old region, this yields the
dimension-independent fresh-coordinate bound in
Lemma~\ref{ap:absolute}. This foundational lifting argument
supplies the absolute rule; preservation of the old width
budgets requires the separate return argument.

\paragraph{Returning to the original cell.}
A density increment found on a sampled parameter box must be
expressed again in the original variables. The cell can have very
small mass, and an error in one of its determining integers
would enter every later polynomial that uses that integer.
The return therefore keeps the old equations exact.

For tests of degree below the current layer weight, we use a
cube version of Conlon--Fox--Zhao densification to replace sparse
factors by bounded ones before applying inverse
theory~\cite{ConlonFoxZhao2015}. The comparison proved here has
bounds independent of that layer's chart mass.
At layers where the testing degree is at least the layer weight,
polynomial families retain the old
equations through an exact projection while allowing the newly
added region to have higher degree. Preparation, ascent and
extraction then control the relative loss of each width by
structural data and higher-weight precision budgets.
The general route from inverse theory to density increments
is classical~\cite{GreenTaoDensity2010}; retaining this dependence
through the return is the quantitative issue addressed here.

\subsection{Coloring thresholds}

Let $W_r(k)$ be the least integer $N$ such that every coloring of
$[N]$ with $r$ colors contains a monochromatic nonconstant
$k$-term arithmetic progression. A largest color class has
density at least $1/r$, so Theorem~\ref{ap:main} gives
\begin{equation}\label{ap:coloring}
 W_r(k)\le
 \left\lceil\exp\bigl(A_k(2+\log r)^{A_k}\bigr)\right\rceil
 \qquad(r\ge1),
\end{equation}
after increasing $A_k$. Thus, for fixed $k$, the sufficient
length is quasipolynomial in the number of colors.
The companion paper~\cite[Theorem~7.2]{OpenAIVanDerWaerden2026}
gives the complementary lower bound
$W_r(k)>\exp((\log r)^2/(64\log2))$ for all integers $r\ge256$
and $k\ge3$. Together these estimates show that for each fixed
$k\ge3$, growth in the number of colors is superpolynomial but
at most quasipolynomial. The upper constant and exponent may
depend on $k$.

\subsection{How to read the proof}

Section~\ref{ap:setup} defines the cells and precision budgets,
states the triangular increment theorem, and explains the
dimension-independent absolute rule.
Figure~\ref{ap:proof-route} gives the main route, and
Table~\ref{ap:interface-table} identifies the analytic statements
and their full proofs.
Section~\ref{ap:paths} constructs the constrained path law and
states its sampling guarantees. Section~\ref{ap:preparation}
prepares the cell by rank cuts, and
Section~\ref{ap:relative-detection} proves the detection bound
independent of the current width. Section~\ref{ap:scalar}
transfers the certificate to terminal boxes and obtains an
absolute increment. Section~\ref{ap:ascent} sets up its return,
develops the comparison estimates, and handles the passive
layers. Section~\ref{ap:active} treats the active layers.
Section~\ref{ap:extraction} extracts the next cell, and
Section~\ref{ap:iteration} closes the recurrences and proves
Theorem~\ref{ap:main}. Section~\ref{ap:weighted} establishes the
weighted consequences.

The appendices contain the rational and symbol calculus, shift
comparison, progression counting, constrained sampling, cube
comparison and relative lifting arguments. The two principal
external analytic inputs to the increment proof are the interval inverse theorem of
Leng, Sah and Sawhney and the radius-sensitive almost-periodicity
theorem of Schoen and Sisask~\cite{LSSInverse2024,SchoenSisask2016}.
Their precise hypotheses are stated at the points of application.
